Generating realistic synthetic option prices requires implied volatility (IV) as an input, yet IV itself is derived from observed option prices, creating a circular dependency that limits synthetic data availability for machine-learning and risk-analysis applications. We broke this circularity by developing a three-stage pipeline in which IV emerged as an output of a structural model of equity returns rather than as an exogenous input. The pipeline coupled a Jump Hidden Markov Model that generated multi-asset price paths with realistic stylized facts and cross-asset tail dependence; a modified Heston stochastic variance process that converted those price paths into implied volatility, where the mean-reversion target was driven by the regime state, days to expiration, moneyness, and an aggregate market mood indicator derived endogenously from the fraction of tickers occupying tail states; and a Cox-Ross-Rubinstein binomial tree that priced American options with early exercise using the resulting IV. A central design choice was initializing the variance process at its mean-reversion target for each strike-expiration pair, so that the smile, skew, and term structure of the IV surface emerged automatically without external calibration data. We calibrated the smile and term-structure shape function through a hierarchy of representations that traded off model capacity for sharing across tickers: a parametric baseline established on a single-ticker multi-expiration chain, a globally shared neural surrogate that removed the functional-form bottleneck, and a sector-specific neural surrogate that let within-sector shapes specialize. We applied this hierarchy to a 31-ticker, six-sector option ladder covering technology, healthcare, financials, energy, retail, and broad-market ETFs, and showed that both added capacity and loosened sharing contributed roughly additive reductions in fit error. Extending the corpus to a multi-day capture and running a temporal holdout further isolated the role of scheduled corporate events: excluding observations near earnings prints on either the contract's own ticker or its same-sector peers eliminated the test-time generalization gap, and reintroducing those observations alongside calendar-derived earnings-distance and peer-distance inputs to the shape network recovered a portion of the lost performance, with the residual gap attributable to surprise-driven post-event moves that no scheduled-event indicator could predict. The same decomposition supports a natural progression to per-ticker neural surrogates as per-name sample sizes grow. The framework was implemented as an open-source Julia package designed to generate synthetic data for downstream machine-learning and risk-analysis applications.
